% !TeX root = ../../book.tex
\section{可数集和不可数集}
\secbegin{secCountableUncountableSets}

在 \Cref{secFiniteSets} 中，我们定义了集合`有限'的含义，以便精确描述元素可以有头有尾一一列出的想法。我们通过观察有限集的元素列表本质上与双射 $f : [n] \to X$ 相同来做到这一点，其中元素 $f(k) \in X$ 扮演列表中第 $k$ 个元素的角色。

我们现在将目光转到\textit{无限}集。我们当然不能期望一个包含无限集中所有元素的列表，所以现在的问题是：如果我们允许列表无限，它的元素是否可以列出？我们将这样的集合称为\textit{可数集}。

令人惊讶的是，并非每个集合都是可数的：有些集合`太大'从而无法列出其元素！我们将在本节后面证明这种\textit{不可数}集合的存在。

集合 $X$ 可数 (\Cref{defCountable}) 的精确定义体现了这样的想法：我们可以逐一列出 $X$ 的元素，这样，即使列表永远继续下去，$X$ 的每个元素都出现在列表中的\textit{在某个有限阶段}。这个列表可能是有限的，也可能不是有限的；它有一个开始，但可能没有结束。

为了说明这一点，请考虑以下 $\mathbb{N}$ 的元素列表：
\[ 0,~ 1,~ 2,~ 3,~ 4,~ \dots,~ n,~ n+1,~ \dots \]
这个列表没有尽头，因为 $\mathbb{N}$ 是无限的（\Cref{thmNIsInfinite}），尽管如此，每个自然数都出现在列表中的某个有限阶段。

另一个例子是所有整数的集合 $\mathbb{Z}$。我们可能希望以通常方式列出 $\mathbb{Z}$ 的元素：
\[ \dots,~ -(n+1),~ -n,~ \dots,~ -3,~ -2,~ -1,~ 0,~ 1,~ 2,~ 3,~ \dots,~ n,~ n+1,~ \dots \]
但这不符合我们的标准，即列表必须有一个开始：它在两个方向上都是无限的。
但是，仍然可以通过将负整数插入非负整数之间来列出它们：
\[ 0,~ -1,~ 1,~ -2,~ 2,~ -3,~ 3,~ \dots,~ -n,~ n,~ -(n+1),~ n+1,~ \dots\]
这不是我们通常列出整数的方式，但我们仍然找到了一种方法，使得每个整数都出现在列表中的某个有限阶段。

但是指定无限集合 $X$ 的元素列表，使得 $X$ 的每个元素出现在列表上的某个有限阶段，相当于指定双射 $f : \mathbb{N} \to X$，其中元素 $f(k) \in X$ 扮演列表中第 $k$ 个元素 --- 也就是说，$X$ 的元素可以列出为
\[  f(0),~ f(1),~ f(2),~ \dots,~ f(n),~ f(n+1),~ \]

这引出了如下定义。

\begin{definition}
\label{defCountable}
\index{countable set}
\index{enumeration!of a countably infinite set}
如果存在双射 $f : \mathbb{N} \to X$，则集合 $X$ 为\textbf{可数无限集}。双射 $f$ 称为 $X$ 的\textbf{枚举}。如果集合 $X$ 是有限集或可数无限集，则说 $X$ 是\textbf{可数的}。
\end{definition}

有些作者更喜欢用`可数'来表示`可数无限'，在这种情况下，他们会说`有限或可数'来表示`可数'。

\begin{example}
\label{exNIsCountable}
集合 $\mathbb{N}$ 为可数无限集，因为根据 \Cref{exIdentityBijection}，恒等函数 $\mathrm{id}_{\mathbb{N}} : \mathbb{N} \to \mathbb{N}$ 为双射。该枚举给出了通常的自然数列表
\[ 0,~ 1,~ 2,~ 3,~ \dots,~ n,~ n+1,~ \dots \]
\end{example}

\begin{example}
\label{exZIsCountable}
函数 $f : \mathbb{Z} \to \mathbb{N}$ 定义为
\[ f(x) = \begin{cases} 2x & \text{如果 } x \ge 0 \\ -(2x+1) & \text{如果 } x < 0 \end{cases} \]
其中 $x \in \mathbb{Z}$，为双射。事实上，它有如下逆
\[ f^{-1}(x) = \begin{cases} \frac{x}{2} & \text{如果 } x \text{ 为偶数} \\ -\frac{x+1}{2} & \text{如果 } x \text{ 为奇数} \end{cases} \]
因此整数集 $\mathbb{Z}$ 是可数无限集。相应的整数列表由下式给出
\[ 0,\ {-1},\ 1,\ {-2},\ 2,\ {-3},\ 3,\ {-4},\ 4,\ \dots \]
这正是我们之前给出的列表！$f$ 是双射的事实确保每个整数在此列表中恰好出现一次。
\end{example}

\begin{exercise}
\label{exPositiveIntegersCountablyInfinite}
证明所有正整数的集合是可数无限集。
\end{exercise}

\begin{exercise}
\label{exEvenOddNaturalNumbersCountablyInfinite}
证明所有偶自然数集合是可数无限集合，且所有奇自然数集合也是可数无限集。
\end{exercise}

由于双射的逆是双射，我们还可以通过找到双射 $X \to \mathbb{N}$ 来证明集合 $X$ 是可数的。

\begin{exercise}
证明由 $p(x,y) = 2^x(2y+1)-1$ 定义的函数 $p : \mathbb{N} \times \mathbb{N} \to \mathbb{N}$ 为双射。推导出 $\mathbb{N} \times \mathbb{N}$ 为可数集。
\hintlabel{exNCrossNIsCountable}{%
算术基本定理 (\Cref{thmFTA}) 告诉我们，对于所有 $n \in \mathbb{N}$，我们可以将 $n+1$ 唯一地表示为 $2$ 的幂乘以某个奇数。
}
\end{exercise}

\subsection*{可数集的封闭性}

通过找到自然数集的显式双射来证明集合是无限集可能过于繁琐，所以我们接下来将发展一些可数集合的\textit{封闭属性}。这样就能够通过将集合 $X$ 与我们已知的可数集合相关联来证明它是可数的，而不必每次都寻找双射 $f : \mathbb{N} \to X$。

\begin{proposition}
\label{propCountablyInfiniteFromBijection}
设 $f : X \to Y$ 为双射。则 $X$ 为可数无限集，当且仅当 $Y$ 为可数无限集。
\end{proposition}

\begin{cproof}
假设 $X$ 为可数无限集。则存在双射 $g : \mathbb{N} \to X$。而 $f \circ g : \mathbb{N} \to Y$ 是双射的组合，所以也为双射，这意味着 $Y$ 是可数无限集。

反过来，假设 $Y$ 为可数无限集。则存在双射 $h : \mathbb{N} \to Y$。而 $f^{-1} \circ h : \mathbb{N} \to X$ 是双射的组合，所以也为双射，这意味着 $X$ 是可数无限集。
\end{cproof}

\begin{exercise}
证明如果 $X$ 和 $Y$ 为可数无限集，则 $X \times Y$ 为可数无限集。
\hintlabel{exProductOfTwoCountableSetsIsCountable}{%
使用 \Cref{exCartesianProductOfBijections} 和可数无限集的定义，来构建双射 $\mathbb{N} \times \mathbb{N} \to X \times Y$。然后应用 \Cref{exNCrossNIsCountable} 和 \Cref{propCountablyInfiniteFromBijection}。
}
\end{exercise}

\Cref{exProductOfTwoCountableSetsIsCountable} 允许我们证明有限多个可数无限集的乘积是可数无限集 --- 这类似可数无限集乘法原理（\Cref{lemMultiplicationPrincipleIndependent}）（的独立版本）。

\begin{proposition}[有限多个可数集合的乘积是可数的]
\label{propFiniteProductOfCountableSetsIsCountable}
设 $n \ge 1$ 并设 $X_1, \dots, X_n$ 为可数无限集。则乘积 $\displaystyle \prod_{i=1}^n X_i$ 为可数无限集。
\end{proposition}
\begin{cproof}
我们在 $n \ge 1$ 上用归纳法来证明。
\begin{itemize}
\item (\textbf{基本情况}) 当 $n=1$ 时，断言不言而喻：如果 $X_1$ 为可数无限集，那么 $X_1$ 为可数无限集。
\item (\textbf{归纳步骤}) 给定 $n \ge 1$ 并假设对于任意集合 $X_1, \dots, X_n$，乘积 $\prod_{i=1}^n X_i$ 为可数无限集。给定集合 $X_1, \dots, X_{n+1}$。那么根据归纳假设 $\prod_{i=1}^n X_i$ 为可数无限集，且根据假设 $X_{n+1}$ 为可数无限集，所以根据 \Cref{exProductOfTwoCountableSetsIsCountable}，集合
\[ \left( \prod_{i=1}^n X_i \right) \times X_{n+1} \]
为可数无限集。根据 \Cref{exProductOfSuccNSets} 存在双射
\[ \prod_{i=1}^{n+1} X_i \to \left( \prod_{i=1}^n X_i \right) \times X_{n+1} \]
所以根据 \Cref{exProductOfTwoCountableSetsIsCountable} 可得 $\prod_{i=1}^{n+1} X_i$ 为可数无限集。
\end{itemize}
利用归纳法，证明完成。
\end{cproof}

我们往往只想知道一个集合是否\textit{可数}，而关心是否是可数无限。这可能是因为我们只需找到集合大小的上限；或者可能是因为我们已经知道集合是无限的，只需证明它是可数的。

下面的定理允许我们通过将 $\mathbb{N}$ 满射到 $X$ 上，或者将 $X$ 单射到 $\mathbb{N}$ 来证明集合 $X$ 是可数的 --- 即有限集或可数无限集。根据 \Cref{secFiniteSets} 的直觉，我们使用单射和满射比较有限集的大小，这表明当且仅当 $X$ 最多与 $\mathbb{N}$ 一样大时，集合才是可数的。

%% Note for future: injection case doesn't require X inhabited.
\begin{theorem} \label{thmCountableFromInjSurj}
设 $X$ 为非空集合。如下陈述是等价的：
\begin{enumerate}[(i)]
\item $X$ 可数；
\item 存在满射 $f : \mathbb{N} \to X$；
\item 存在单射 $f : X \to \mathbb{N}$。
\end{enumerate}
\end{theorem}
\begin{cproof}
我们来证明 (i)$\Leftrightarrow$(ii) 和 (i)$\Leftrightarrow$(iii)。
\begin{itemize}
\item (i)$\Rightarrow$(ii). 假设 $X$ 可数。如果 $X$ 为可数无限集，则存在双射 $f : \mathbb{N} \to X$，也是满射。如果 $X$ 为有限集，则存在双射 $g : [m] \to X$，其中 $m = |X| \ge 1$。定义函数 $f : \mathbb{N} \to X$ 为
\[ f(n) = \begin{cases} g(n) & \text{ 如果 } 1 \le n \le m \\ g(1) & \text{ 如果 } n=0 \text{ 或 } n > m \end{cases} \]
则 $f$ 为满射：如果 $x \in X$ 则存在 $n \in [m]$ 使得 $g(n)=x$，那么 $f(n)=g(n)=x$。
\item (ii)$\Rightarrow$(i). 假设存在满射 $f : \mathbb{N} \to X$。要证明 $X$ 可数，只要证明如果 $X$ 是无限集则它是可数无限集即可。所以假设 $X$ 是无限集，并递归地定义序列
\begin{itemize}
\item $a_0=0$;
\item 给定 $n \in \mathbb{N}$ 并假设 $a_0, \dots, a_n$ 已经定义。定义 $a_{n+1}$ 为 $f(a_{n+1}) \not \in \{ f(a_0), f(a_1), \dots, f(a_n) \}$ 的最小自然数。
\end{itemize}
对于所有 $n \in \mathbb{N}$ 定义 $g : \mathbb{N} \to X$ 为 $g(n)=f(a_n)$。则
\begin{itemize}
\item $g$ 是单射的。因为如果 $m \le n$ 则 $f(a_m) \ne f(a_n)$ by construction of the sequence $(a_n)_{n \in \mathbb{N}}$.
\item $g$ 是满射的。事实上，给定 $x \in X$，根据满射性，存在最小 $m \in \mathbb{N}$ 使得 $f(m)=x$，因此对于某个 $n \le m$ 通过序列 $(a_n)_{n \in \mathbb{N}}$ 必然可得 $a_n=m$。所以 $x = f(a_n) = g(n)$，因此 $g$ 是满射的。
\end{itemize}
所以 $g$ 为双射，且 $X$ 是可数的。
\item (i)$\Rightarrow$(iii). 假设 $X$ 是可数的。如果 $X$ 为可数无限集，则存在双射 $f : \mathbb{N} \to X$，所以 $f^{-1} : X \to \mathbb{N}$ 也是双射的，因此也是单射的。如果 $X$ 为有限集，则存在双射 $g : [m] \to X$，其中 $m=|X| \ge 1$。那么 $g^{-1} : X \to [m]$ 是单射的。设对于所有 $k \in [m]$, $i : [m] \to \mathbb{N}$ 定义为 $i(k)=k$。则 $i \circ g^{-1}$ 是单射的；事实上, 对于 $x,x' \in X$ 我们有
\[ i(g^{-1}(x))=i(g^{-1}(x')) \Rightarrow g^{-1}(x)=g^{-1}(x') \Rightarrow x=x' \]
第一点由 $i$ 的定义得到，第二点由 $g^{-1}$ 的单射性得到。所以存在单射 $X \to \mathbb{N}$。
\item (iii)$\Rightarrow$(i). 假设存在单射 $f : X \to \mathbb{N}$。要证明 $X$ 为可数集，只需证明如果 $X$ 是无限集则它为可数无限集即可。递归地定义序列 $(a_n)_{n \in \mathbb{N}}$ 如下：
\begin{itemize}
\item 设 $a_0$ 为 $f[X]$ 的最小元素；
\item 给定 $n \in \mathbb{N}$ 并假设 $a_0, \dots, a_n$ 已经定义。设 $a_{n+1}$ 为 $f[X] \setminus \{ a_0, \dots, a_n \}$ 的最小元素。因为 $f$ 为单射，所以 $f[X]$ 是无限的。 
\end{itemize}
定义 $g : \mathbb{N} \to X$ ，对于每个 $n \in \mathbb{N}$，设 $g(n)$ 是 $x$ 的唯一值，其中 $f(x)=a_n$。那么
\begin{itemize}
\item $g$ 是单射的。每当 $m \ne n$，构造 $a_m \ne a_n$。设 $x,y \in X$ 使得 $f(x)=a_m$ 且 $f(y)=a_n$。因为 $f$ 是单射的，可得 $x \ne y$，所以 $g(m) = x \ne y = g(n)$。
\item $g$ 是满射的。给定 $x \in X$。那么 $f(x) \in f[X]$，所以存在 $m \in \mathbb{N}$ 使得 $f(x)=a_m$。因此 $g(m)=x$。
\end{itemize}
所以 $g$ 为双射，且 $X$ 为可数无限集。
\end{itemize}
以上证明了上述等价性。
\end{cproof}

实际上，我们不需要用 $\mathbb{N}$ 作为满射的定义域或单射的值域；事实上我们可以使用任意可数集 $C$。

\begin{exercise}
\label{exCountabilityFromInjectionsAndSurjections}
设 $X$ 为非空即可。如下陈述是等价的：
\begin{enumerate}[(i)]
\item \label{itmInjectionSurjectionCountableSets} $X$ 是可数的；
\item \label{itmInjectionIntoCountableSet} 对于可数集$C$，存在单射 $f : X \to C$；
\item \label{itmSurjectionFromCountableSet} 对于可数集$C$，存在满射 $f : C \to X$。
\end{enumerate}
\hintlabel{exCountableFromInjSurjC}{%
对于 \ref{itmInjectionIntoCountableSet}$\Rightarrow$\ref{itmInjectionSurjectionCountableSets}，利用两个单射的复合为单射这一性质。对于 \ref{itmSurjectionFromCountableSet}$\Rightarrow$\ref{itmInjectionSurjectionCountableSets} 也是如此。
}
\end{exercise}

\Cref{exCountableFromInjSurjC} 对证明许多其他集合的可数性非常有用：当我们建立可数集的全部技能时，为了证明集合 $X$ 是可数的，我们需要做的是找到从已知可数集到 $X$ 的满射，或从 $X$ 到已知可数集的单射。

这一证明技巧为下面反直觉的结论给出了令人难以置信的简短证明，这个反直觉的结论可以解释为有理数的数量与自然数的数量完全相等。

\begin{theorem}
\label{thmRationalsAreCountable}
有理数集 $\mathbb{Q}$ 是可数的。
\end{theorem}

\begin{cproof}
定义函数 $q : \mathbb{Z} \times (\mathbb{Z} \setminus \{ 0 \}) \to \mathbb{Q}$，对于所有 $a,b \in \mathbb{Z}$ 且 $b \ne 0$，设 $q(a,b) = \frac{a}{b}$。

根据 \Cref{exZIsCountable} 和 \Cref{exProductOfTwoCountableSetsIsCountable}，集合 $\mathbb{Z} \times (\mathbb{Z} \setminus \{ 0 \})$ 是可数的。

根据 $\mathbb{Q}$ 的定义，函数 $q$ 是满射的 --- 事实上，$x \in \mathbb{Q}$ 的精确表达是 $x = \dfrac{a}{b} = q(a,b)$ 其中 $(a,b) \in \mathbb{Z} \times (\mathbb{Z} \setminus \{ 0 \})$。

根据 \Cref{exCountableFromInjSurjC}，可得 $\mathbb{Q}$ 是可数的。
\end{cproof}

\begin{exercise} \label{exFiniteSubsetsCountableFixedSize}
设 $X$ 为可数集。证明对于每个 $k \in \mathbb{N}$, $\binom{X}{k}$ 是可数的。
\begin{backhint}
\hintref{exFiniteSubsetsCountableFixedSize}
假设 $X=\mathbb{N}$。根据 \Cref{propFiniteProductOfCountableSetsIsCountable}，集合 $\mathbb{N}^k$ 是可数的。根据 \Cref{thmCountableFromInjSurj}(c)，只需找到单射 $\binom{\mathbb{N}}{k} \to \mathbb{N}^k$ 即可。
\end{backhint}
\end{exercise}

\begin{theorem}[可数多个可数集得并集是可数的]
\label{thmCountableUnionOfCountableSetIsCountable}
设 $\{ X_n \mid n \in \mathbb{N} \}$ 是一个可数集族。则集合 $X$ 
\[ X = \bigcup_{n \in \mathbb{N}} X_n \]
是可数的。
\end{theorem}

\begin{cproof}
我们假设集合 $X_n$ 都是非空的，因为空集对并运算没贡献。

对于每个 $n \in \mathbb{N}$ 存在满射 $f_n : \mathbb{N} \to X_n$。定义函数 $f : \mathbb{N} \times \mathbb{N} \to X$ 为 $f(m,n)=f_m(n)$ 其中 $m,n \in \mathbb{N}$。则 $f$ 是满射的：如果 $x \in X$ 则 $x \in X_m$ 其中 $m \in \mathbb{N}$。因为 $f_m$ 是单射的，可得 $x=f_m(n)$ 其中 $n \in \mathbb{N}$。那么 $x=f(m,n)$。因为 $\mathbb{N} \times \mathbb{N}$ 是可数的，根据 \Cref{exCountableFromInjSurjC} 可得 $X$ 是可数的。
\end{cproof}

\begin{example}
\label{exCountableSubsetsOfCountableSetIsCountable}
设 $X$ 为可数集。$X$ 的所有有限子集的集合是可数的。事实上，$X$ 的所有有限子集的集合等于 $\displaystyle \bigcup_{k \in \mathbb{N}} \binom{X}{k}$，根据 \Cref{exFiniteSubsetsCountableFixedSize}，它是可数多个可数集的并集，所以根据 \Cref{thmCountableUnionOfCountableSetIsCountable} 它是可数的。
\end{example}

\subsection*{不可数集}

我们已经见过太多可数集的例子了。并非每个集合都是可数的，这一点并不是显而易见的。我们如何\textit{知道}存在某些集合，其元素无法列出呢？

在 \Cref{thmCantorForN} 中，我们证明了存在一个不可数集，即自然数的幂集。证明看似简单，但其意义非常重要。

\begin{theorem}
\label{thmCantorForN}
$\mathcal{P}(\mathbb{N})$ 是不可数的。
\end{theorem}

\begin{cproof}
在 \Cref{exRussellSubset} 中，我们证明了没有函数 $f : \mathbb{N} \to \mathcal{P}(\mathbb{N})$ 是满射的。具体来说，给定 $f : \mathbb{N} \to \mathcal{P}(\mathbb{N})$，我们可以证明对于任意 $n \in \mathbb{N}$，元素
\[ B = \{ k \in \mathbb{N} \mid k \not\in f(k) \} \in \mathcal{P}(\mathbb{N}) \]
不等于 $f(n)$。 根据 \Cref{exCountabilityFromInjectionsAndSurjections}\ref{itmSurjectionFromCountableSet} 可知 $\mathcal{P}(\mathbb{N})$ 是不可数的。
\end{cproof}

\Cref{thmCantorForN} 中使用的方法是\textit{康托尔对角线法}的一个典型示例。

\begin{theorem}
\label{thmCantorDiagonalGeneral}
设 $X$ 为集合，并假设对于每个函数 $f : \mathbb{N} \to X$，都有:
\begin{enumerate}[(i)]
\item 逻辑公式族 $p_n(x)$ 其中 $x \in X$，每个 $n \in \mathbb{N}$ 一个；
\item 元素 $b \in X$；
\end{enumerate}
满足 $\forall n \in \mathbb{N},~ [p_n(b) \Leftrightarrow \neg p_n(f(n))]$。则 $X$ 是不可数的。
\end{theorem}

\begin{cproof}
我们证明没有函数 $f : \mathbb{N} \to X$ 是满射的。所以设 $f : \mathbb{N} \to X$ 为任意函数，并设 $p_n(x)$ 和 $b \in X$ 与定理的陈述相同。

要想证明 $f$ 不是满射的, 先假设 $b=f(k)$ 其中 $k \in \mathbb{N}$，然后引出矛盾。
\begin{itemize}
\item 如果 $p_k(b)$ 为真，则 $p_k(f(k))$ 为真，因为 $b=f(k)$，且 $\neg p_k(f(k))$ 为真，因为 $p_k(b) \Rightarrow \neg p_k(f(k))$。存在矛盾。
\item 如果 $\neg p_k(b)$ 为真，则 $\neg p_k(f(k))$ 为假，因为 $b=f(k)$，所以 $p_k(b)$ 为真，因为 $[\neg p_k(f(k))] \Rightarrow p_k(b)$。存在矛盾。
\end{itemize}
在这两种情况下，我们都得到矛盾。所以对于任意 $k \in \mathbb{N}$, $b \ne f(k)$，因此 $f$ 不是满射的。

因此，根据 \Cref{exCountabilityFromInjectionsAndSurjections}\ref{itmSurjectionFromCountableSet}，$X$ 是不可数的。
\end{cproof}

\begin{strategy}[康托尔对角线法]
\label{strCantorDiagonal}
\index{Cantor's diagonal argument}
要证明集合 $X$ 是不可数的，只需用如下方法证明没有函数 $f : \mathbb{N} \to X$ 是满射的：给定函数 $f : \mathbb{N} \to X$，找到元素 $b \in X$ `不满足'涉及 $n$ 的某个陈述的每个值 $f(n)$。
\end{strategy}

\begin{example}
在 \Cref{thmCantorForN} 的证明中，我们证明了 $\mathcal{P}(\mathbb{N})$ 是不可数的。方法如下，给定 $f : \mathbb{N} \to \mathcal{P}(\mathbb{N})$，我们需要找到元素 $B \in \mathcal{P}(\mathbb{N})$ --- 也就是说，子集 $B \subseteq \mathbb{N}$ --- `不满足'涉及 $n$ 的每个值 $f(n)$。

记住，对于每个 $n \in \mathbb{N}$，元素 $f(n) \in \mathcal{P}(\mathbb{N})$ 是 $\mathbb{N}$ 的子集，因此询问 $n$ 是否是 $f(n)$ 的元素是有意义的。

考虑到这一点，对于每个 $n \in \mathbb{N}$，我们迫使 $B$ 在 $n$ 是否为元素的问题上与 $f(n)$ 不一致。也就是
\[ n \in B \quad \Leftrightarrow \quad n \not\in f(n) \]
用 \Cref{thmCantorDiagonalGeneral} 的语言，陈述 $p_n(U)$ （对于 $n \in \mathbb{N}$ 和 $U \in \mathcal{P}(\mathbb{N})$）为陈述 `$n \in U$'；对于每个 $n \in \mathbb{N}$，这个陈述在 $U=B$ 时为当且仅当其在 $U=f(n)$ 时为假。

$B = \{ k \in \mathbb{N} \mid k \not\in f(k) \}$ 的定义迫使对于所有 $n \in \mathbb{N}$, $n \in B \Leftrightarrow n \not\in f(n)$ 为真，因此康托尔对角线法适用。
\end{example}

康托尔对角线法还可以通过实数的十进制展开（\Cref{defBaseBExpansionOfRealNumber}）来证明 $\mathbb{R}$ 是不可数的。

\begin{theorem}
\label{thmRIsUncountable}
$\mathbb{R}$ 是不可数的。
\end{theorem}

\begin{cproof}
设 $f : \mathbb{N} \to \mathbb{R}$ 为任意函数。我们证明 $f$ 不是满射。

对于每个 $n \in \mathbb{N}$，设 $f(n) \in \mathbb{R}$ 的十进制展开为 $x_0^n \,.\, x_1^n x_2^n x_3^n \dots$ --- 也就是说
\[ f(n) = \sum_{i \ge 0} x_i^n \cdot 10^{-i} \]
[这里的上标 `$n$' 是一个标签，而不是指数。]

我们通过强制 $b$ 和 $f(n)$ 的十进制展开在每个 $n \in \mathbb{N}$ 上的第 $n$ 位不同来定义 $b \in \mathbb{R}$。具体来说，对于每个 $n \in \mathbb{N}$ 设
\[ b_n = \begin{cases} 0 & \text{如果 } x_n^n \ne 0 \\ 1 & \text{如果 } x_n^n = 0 \end{cases} \]
并设 $b = b_0. b_1 b_2 b_3 \dots{} = \sum_{i \ge 0} b_i \cdot 10^{-i}$。

请注意，$b$ 没有重复出现的9，因此它是实数的有效十进制展开。并且对于任意 $n \in \mathbb{N}$ 根据实数十进制展开的唯一性（\Cref{thmBaseBExpansionsOfReals}），$b \ne f(n)$，因为 $b_n \ne x_n^n$。

因此 $f$ 不是满射的，所以 $\mathbb{R}$ 是不可数的。
\end{cproof}

\begin{exercise}
用康托尔对角线法证明所有函数集合 $\mathbb{N} \to \mathbb{N}$ 的集合 $\mathbb{N}^{\mathbb{N}}$ 是不可数的。
\end{exercise}

每次我们想要证明一个集合不可数时，使用康托尔对角线法可能会很麻烦。然而，就像 \Cref{thmCountableFromInjSurj} 一样，我们可以通过将集合与我们已知的其他不可数集相关联来证明集合是不可数的。

\begin{theorem}
\label{thmUncountableFromInjSurj}
设 $X$ 为集合。如下陈述是等价的：
\begin{enumerate}[(i)]
\item $X$ 不可数；
\item \label{testitem} 对于不可数集 $U$，存在满射 $f : X \to U$；
\item 对于不可数集 $U$，存在单射 $f : U \to X$。
\end{enumerate}
\end{theorem}

\begin{cproof}
对于 (i)$\Rightarrow$(ii) 和 (i)$\Rightarrow$(iii)，令 $U=X$ 且 $f=\mathrm{id}_X$。

对于 (ii)$\Rightarrow$(i)，对于不可数集 $U$，假设存在满射 $f : X \to U$。如果 $X$ 可数，则根据 \Cref{thmCountableFromInjSurj}(ii)，存在满射 $g : \mathbb{N} \to X$；那么 $f \circ g : \mathbb{N} \to U$ 为满射，这与 $U$ 不可数矛盾。

对于 (iii)$\Rightarrow$(i)，对于不可数集 $U$，假设存在单射 $f : U \to X$。如果 $X$ 可数，则根据 \Cref{thmCountableFromInjSurj}(iii)，存在单射 $g : X \to \mathbb{N}$；那么 $g \circ f : U \to \mathbb{N}$ 为单射, 这与 $U$ 不可数矛盾。
\end{cproof}

\begin{proposition}
\label{propBinarySequencesAreUncountable}
所有函数 $\mathbb{N} \to \{0,1\}$ 的集合 $\{0,1\}^{\mathbb{N}}$ 是不可数的。
\end{proposition}

\begin{cproof}
定义函数 $F : \mathcal{P}(\mathbb{N}) \to \{0,1\}^{\mathbb{N}}$，对于 $U \subseteq \mathbb{N}$ 令 $F(U) = \chi_U : \mathbb{N} \to \{0,1\}$ 为 $U$ 的特征函数（\Cref{defCharacteristicFunction}）。

给定 $U, V \subseteq \mathbb{N}$，根据 by \Cref{thmCharacteristicFunctionsCharacteriseSubsets}，我们有
\[ F(U) = F(V) \quad \Rightarrow \quad \chi_U = \chi_V \quad \Rightarrow \quad U = V \]
则 $F$ 是单射的，而根据康托尔定理 $\mathcal{P}(\mathbb{N})$ 是不可数的，所以根据 \Cref{thmUncountableFromInjSurj}，$\{0,1\}^{\mathbb{N}}$ 是不可数的。
\end{cproof}

\begin{exercise}
证明如果集合 $X$ 拥有不可数子集，则 $X$ 是不可数的。
\end{exercise}

\begin{exercise}
\label{exPowerSetFiniteOrUncountable}
设 $X$ 为集合。证明 $\mathcal{P}(X)$ 要么是有限的，要么是不可数的。
\begin{backhint}
\hintref{exPowerSetFiniteOrUncountable}
我们已经证明 $X$ 为有限集的情况。当 $X$ 为可数无限集时，用 \Cref{thmCantorForN}。当 $X$ 为不可数无限集时，找到单射 $X \to \mathcal{P}(X)$ 并想办法应用 \Cref{exCountableFromInjSurjC}。
\end{backhint}
\end{exercise}

\subsection*{检测可数性}

我们已经见到了若干可数集和不可数集的例子。但是，能够证明一个集合是可数的，或者一个集合是不可数的，技只此耳 --- 能够获得如何检测集合可数性的直觉大有益处。

我们接下来致力于证明以下启发：

\begin{center}
\textit{如果集合的每个元素都有有限的描述，则该集合是可数的。}
\end{center}

在我们进行精确阐述之前，让我们看一下这个启发式的实际应用。考虑如下我们已经见过的可数集合：
\begin{itemize}
\item 集合 $\mathbb{N}$ 是可数的；每个自然数都可以用一个有限的数字字符串表示，因此存在一个有限的描述 --- 例如，$7$ 或 $15213$。
\item 集合 $\mathbb{Z}$ 是可数的；每个整数都可以用一个有限的数字字符串表示，可能会带上负号 --- 例如，$7$ 或 $-15213$。
\item 集合 $\mathbb{Q}$ 是可数的；每个有理数都都可以用有限个整数（整数本身具有有限的描述）对来表示，一个在另一个之上，用横线分割 --- 例如，$\frac{7}{1}$ 或 $\frac{-15213}{21128}$。
\item $\mathbb{N}$ 的所有\textit{有限}子集的集合是可数的；$\mathbb{N}$ 的每个有限子集在列表表示法中都有一个有限的表示，一个开大括号`$\{$'，后跟有限个自然数（自然数本身具有有限的表达式）以逗号分割，最后写一个闭大括号`$\}$' --- 例如 $\{ 15, 0, 281 \}$ 或 $\{ \}$。
\end{itemize}

接下来考虑我们已知（或被告知）的不可数集：
\begin{itemize}
\item 集合 $\mathcal{P}(\mathbb{N})$ 是不可数的；直观上，$\mathbb{N}$ 的无限子集不能以统一的方式有限地表示。 
\item 集合 $\mathbb{R}$ 是不可数的；直观上，没有统一的方法来有限地描述实数 --- 例如，十进制展开式就不够好，因为它们是无限的。
\end{itemize}

\begin{definition}
\label{defKleeneStar}
\index{Kleene star}
\index{word}
\index{empty word}
\index{length!of a word}
\index{alphabet}
\nindex{Sigmastar}{$\Sigma^*$}{Kleene star}
设 $\Sigma$ 为集合。$\Sigma$ 上的\textbf{单词}是 $a_1 a_2 \dots a_n$ 形式的表达式，其中 $n \in \mathbb{N}$ 且对于所有 $i \in [n]$, $a_i \in \Sigma$。自然数 $n$ 被称为单词的\textbf{长度}。长度为 $0$ 的唯一单词被称为\textbf{空词}，并用 $\varepsilon$ \inlatex{varepsilon} 表示。

集合 $\Sigma$ 被称为\textbf{字母表}。$\Sigma$ 上所有单词的集合用 $\Sigma^*$ \inlatex{Sigma\^{}*} 表示；运算 $(-)^{\star}$ 被称为 \textbf{克莱恩星}.
\end{definition}

形式上，$\Sigma^*$ 被定义为并集 $\displaystyle\bigcup_{n \in \mathbb{N}} \Sigma^n$，其中集合 $\Sigma^n$ 递归地定义为
\[ \Sigma^0 = \{ \varepsilon \} \quad \text{且} \quad \Sigma^{n+1} = \{ wa \mid w \in \Sigma^n,~ a \in \Sigma \} \text{ 对于所有 } n \ge 0 \]
即，给定 $n \in \mathbb{N}$，$\Sigma^n$ 的元素正是 $\Sigma$ 上长度为 $n$ 的单词。

\begin{example}
\label{exSomeExamplesOfFiniteStrings}
$\{ 0,1 \}^*$ 的元素正是有限二进制串：
\[ \{ 0,1 \}^* = \{ \varepsilon,~ 0,~ 1,~ 00,~ 01,~ 10,~ 11,~ 000,~ 001,~ 010,~ 011,~ \dots \} \]
同理，$\{ 0,1,2,3,4,5,6,7,8,9 \}^*$ 的元素是以 $10$ 为基数的有限数字串。但请牢记，正式地讲这些串的元素是\textit{单词}而非\textit{数字}，并且两个单词可能对应于相同的自然数 --- 例如 $123$ 和 $00123$ 是 $\{ 0,1,2,3,4,5,6,7,8,9 \}^*$ 的不同元素，尽管事实上它们\textit{代表}相同的自然数。
\end{example}

\begin{exercise}
设 $\Sigma = \{ 0,1,2,3,4,5,6,7,8,9,{\div},{-} \}$。描述 $\Sigma$ 上的哪些单词代表有理数（将数字符号视为数字并将符号 $\div$ 和 $-$ 视为算术运算）。找出 $\Sigma^*$ 的元素无法代表有理数的例子。
\hintlabel{exWordsRepresentingRationalNumbers}{%
如果 $\Sigma^*$ 中的字符串要表示有理数，那么该字符串可以有多少个`$\div$'符号？$\Sigma$ 上的单词中在什么位置可以出现 $\div$ 符号？
}
\end{exercise}

\begin{exercise}
设 $\Sigma$ 为字母表。证明如果 $\Sigma$ 是可数的，则 $\Sigma^*$ 是可数的。
\hintlabel{exIfAlphabetCountableThenSetOfWordsCountable}{%
用归纳法证明对于所有 $n \in \mathbb{N}$, $\Sigma^n$ 是可数的，然后应用 \Cref{thmCountableUnionOfCountableSetIsCountable}。
}
\end{exercise}

正如 \Cref{exSomeExamplesOfFiniteStrings} 暗示的那样，我们可以使用集合上的单词来精确表达`有限描述'的含义。

其思想是，有限地描述集合 $X$ 的元素相当于对于某集合 $\Sigma$ 定义一个函数 $D : X \to \Sigma^*$；给定 $x \in X$，单词 $D(x) \in \Sigma^*$ 可以被认为是元素 $x \in X$ 的有限描述。 

但 $D$ 不能是任何旧函数 --- 如果 $X$ 的两个元素具有相同的描述，则他们应该是相等的。即
\[ \forall x,y \in X,~ D(x) = D(y) \Rightarrow x=y \]
而这恰恰表明 $D$ 必须是单射的！

\begin{definition}
\label{defFiniteDescription}
\index{finite description}
设 $X$ 为集合，并设 $\Sigma$ 为字母表。$X$ 中的元素在 $\Sigma$ 上的\textbf{有限描述}为单射 $D : X \to \Sigma^*$。
\end{definition}

\begin{example}
\label{exFiniteDescriptionOfFiniteSubsetsOfN}
设 $\mathcal{F}(\mathbb{N})$ 为 $\mathbb{N}$ 的所有有限子集的集合。根据 \Cref{exCountableSubsetsOfCountableSetIsCountable}，我们已知此集合是可数的；现在我们演示其元素的有限描述。

设 $\Sigma = \left\{ 0,1,2,3,4,5,6,7,8,9,\boxed{\{},\boxed{\}},\boxed{,} \right\}$ --- 即，字母表 $\Sigma$ 的元素为数字 $0$ 到 $9$，左大括号，右大括号和逗号。 （大括号和逗号外面加了框，以区别列表表示法表示 $\Sigma$ 用到的相同符号。）

定义 $D : \mathcal{F}(\mathbb{N}) \to \Sigma^*$，令 $D(U)$ 为有限集 $U \subseteq \mathbb{N}$ 的列表表示，其中其元素按升序列出。这确保了 $D$ 的准确定义--- 例如
\[ \{1,2,3\} = \{2,3,1\} \quad \text{而} \quad D(\{1,2,3\}) = D(\{2,3,1\}) = \{1,2,3\} \] 

那么 $D$ 是单射的：给定有限集 $U, V \subseteq \mathbb{N}$，如果 $D(U) = D(V)$，则 $U$ 和 $V$ 在列表表示法中具有相同的表达式 --- 这意味着它们具有相同的元素，因此它们是相等的！
\end{example}

\begin{example}
\label{exTrivialFiniteDescription}
每个集合都有一个简单的有限描述，即将集合本身视为字母表。事实上，对于所有 $x \in X$，由 $D(x) = x$ 定义的函数 $D : X \to X^*$ 显然是单射。
\end{example}

\Cref{exTrivialFiniteDescription} 提出了这样一个观点：如果我们想通过证明集合元素具有有限描述来证明集合是可数的，那么我们必须对可能使用的字母表施加一些限制。例如，我们在 \Cref{exFiniteDescriptionOfFiniteSubsetsOfN} 中使用的字母表是有限的 --- 它有 $13$ 个元素。

\begin{exercise}
\label{exFiniteDescriptionOfQ}
在有限字母表上给出有理数的显式有限描述。
\end{exercise}

\begin{exercise}
\label{exFiniteDescriptionOfCofiniteSubsetsOfN}
设 $\mathcal{C}(\mathbb{N})$ 为 $\mathbb{N}$ 的所有\textit{余有限}子集的集合；即
\[ \mathcal{C}(\mathbb{N}) = \{ U \subseteq \mathbb{N} \mid \mathbb{N} \setminus U \text{ 为有限集} \} \]
在有限字母表上给出 $\mathcal{C}(\mathbb{N})$ 元素的显式有限描述。
\end{exercise}

\begin{theorem}
\label{thmCountableIffFiniteDescription}
设 $X$ 为集合。则 $X$ 是可数的，当且仅当 $X$ 的元素在可数字母表上存在有限描述。
\end{theorem}

\begin{cproof}
如果 $X$ 是可数的，则我们可以取 $X$ 本身作为其字母表！对于所有 $x \in X$，由 $D(x) = x$ 定义的函数 $D : X \to X^*$ 显然是单射。

反过来，对于某可数字母表 $\Sigma$，如果存在单射 $D : X \to \Sigma^*$，则根据 \Cref{exIfAlphabetCountableThenSetOfWordsCountable}，$\Sigma^*$ 是可数的，因此根据 \Cref{exCountabilityFromInjectionsAndSurjections}\ref{itmInjectionIntoCountableSet}，$X$ 是可数的。
\end{cproof}

\begin{example}
根据 \Cref{thmCountableIffFiniteDescription}，我们可以从 \Cref{exFiniteDescriptionOfFiniteSubsetsOfN} 和 \Cref{exFiniteDescriptionOfQ,exFiniteDescriptionOfCofiniteSubsetsOfN} 推导出 $\mathcal{F}(\mathbb{N})$, $\mathbb{Q}$ and $\mathcal{C}(\mathbb{N})$ 都是可数集。
\end{example}

\begin{exercise}
证明所有具有有理系数的多项式的集合是可数的。
\end{exercise}

\begin{exercise}
考虑以下论证，该论证声称证明了 $\mathcal{P}(\mathbb{N})$ 是可数的。（事实并非如此。）
\begin{quote}
设 $\Sigma = \left\{ \mathbb{N},x,U,V,\boxed{\mid},\boxed{\in},\boxed{\{},\boxed{\}} \right\}$ --- 这里我们再一次用框框起某些符号来强调它们是 $\Sigma$ 的元素。

对于所有 $U, V \subseteq \mathbb{N}$，定义 $D : \mathcal{P}(\mathbb{N}) \to \Sigma^*$ 为
\[ D(U) = \{ x \in \mathbb{N} \mid x \in U \} \quad \text{且} \quad D(V) = \{ x \in \mathbb{N} \mid x \in V \} \]
则 $\Sigma$ 是有限的，且 $D$ 为单射，因为对于所有 $U,V \subseteq \mathbb{N}$ 我们有
\[ D(U) = D(V) ~ \Rightarrow ~ \{ x \in \mathbb{N} \mid x \in U \} = \{ x \in \mathbb{N} \mid x \in V \} ~ \Rightarrow ~ U = V \]
因为 $D$ 为 $\mathbb{N}$ 的子集在有限字母表上的有限描述，根据 \Cref{thmCountableIffFiniteDescription} 可得 $\mathcal{P}(\mathbb{N})$ 是可数的。\hfill $\Box$
\end{quote}
这个证明是不正确的，因为我们知道 $\mathcal{P}(\mathbb{N})$ 是不可数的。证明中哪里出错了？
\hintlabel{exBadProofThatPowerSetOfNIsCountable}{%
$\Sigma$ \textit{真的}是有限的吗？$D$ \textit{真的}是明确定义的吗？
}
\end{exercise}

\begin{exercise}
假设我们用与 \Cref{exFiniteDescriptionOfFiniteSubsetsOfN} 相同的方法尝试证明 $\mathcal{P}(\mathbb{N})$ 是可数的。
证明会在哪里失败？
\end{exercise}

\begin{exercise}
证明，对于每个可数集合 $X$, $X$ 的元素在\textit{有限}字母表上存在有限描述。
\hintlabel{exCountableIffFiniteDescriptionFiniteAlphabet}{%
首先证明在有限字母表 $\Phi$ 上 $\mathbb{N}^*$ 的元素存在有限描述。这定义了单射 $\mathbb{N}^* \to \Phi^*$。接着给定可数字母表 $\Sigma$，使用 \Cref{thmCountableFromInjSurj} 给出的单射 $\Sigma \to \mathbb{N}$ 构建单射 $\Sigma^* \to \mathbb{N}^*$，从而得到单射 $\Sigma^* \to \Phi^*$。最后，用它来将 $\Sigma$ 上的有限描述转换为 $\Phi$ 上的有限描述。
}
\end{exercise}